\section*{Chapitre \ref{ch:nw:futures-connectivity} --- Connectivité des contrats à terme}

\begin{solution}{exo:nw:futures-connectivity:1}
Une session Convenience Gateway atteint chaque segment de marché, la passerelle acheminant chaque message~; une connexion à une \omterm{def:nw:futures-connectivity:msgw}{passerelle de segment de marché}
n'atteint qu'un seul segment, un groupe de produits sur un même moteur, et la firme se connecte à la passerelle de chaque segment avec la même
identité de session.
\end{solution}

\begin{solution}{exo:nw:futures-connectivity:2}
$6 \times 260 + 4 \times 520 = 3\,640$~euros par mois, pour $10 \times 150 = 1\,500$ transactions par seconde.
\end{solution}

\begin{solution}{exo:nw:futures-connectivity:3}
Un paquet manqué sur un flux se retrouve en général sur l'autre, si bien que traiter les deux évite la plupart des lacunes~; une lacune sur les deux signifie que chaque carnet d'ordres
du canal doit être resynchronisé, ce qui coûte du temps précisément quand les données sont les plus nombreuses.
\end{solution}

\begin{solution}{exo:nw:futures-connectivity:4}
Toutes les copies de l'ordre de chaque firme atteignent ensemble la carte de l'unique passerelle, et sa file est premier entré, premier sorti~: la première copie est
servie au plus tôt au même moment qu'une copie unique le serait, si bien que le gagnant est la firme dont l'ordre arrive en premier, à ±2.5~ns de câble près,
tiré au sort uniformément parmi dix firmes également rapides.
\end{solution}

\begin{solution}{exo:nw:futures-connectivity:5}
Envoyer le début d'un message avant que sa fin soit prête permet à une firme de commencer sa transmission plus tôt, un moyen d'arriver en premier que la
conception ne visait pas à récompenser. La session est retardée d'au moins \qty{3}{\micro\second} et un message complet provenant d'une autre session
peut atteindre le moteur en premier.
\end{solution}

\begin{solution}{exo:nw:futures-connectivity:6}
Que les quatre puissent être exécutées~: ses limites de position doivent compter quatre fois la taille de l'ordre tant que les autres copies ne sont pas annulées ou
n'expirent, et sa réconciliation doit prévoir jusqu'à quatre exécutions.
\end{solution}

\begin{solution}{exo:nw:futures-connectivity:7}
Cinq sessions Full et neuf Ultra, 8\,320~euros par mois~; vingt sessions Full coûtent 8\,840~euros, soit 520~euros de plus, parce que quatorze d'entre elles
sont au prix applicable à partir de la septième.
\end{solution}

\begin{solution}{exo:nw:futures-connectivity:8}
Si les autres copient, le gain disparaît (un sur dix de nouveau) pendant que la latence de chacun augmente, et ajouter des sessions ne fait que relancer la
course, à un coût plus élevé~; et la bourse, dont le système porte la charge, a toutes les raisons de changer la conception, comme l'a fait Eurex.
\end{solution}

\begin{solution}{pb:nw:futures-connectivity:1}
\textbf{Partie I.}
\begin{enumerate}
\item 10.0~\% et 59.7~\%.
\item 8.4 points de pourcentage (840~euros par mois avec les valeurs du chapitre), et 1.8 points de pourcentage (180.5~euros).
\item Le minimum des délais indépendants diminue à chaque copie (la \cref{prop:nw:futures-connectivity:copies}).
\item Les files d'attente indépendantes des passerelles et la gigue (jitter), par rapport aux écarts entre firmes~; sans gigue, il n'y aurait rien
  à gagner.
\end{enumerate}
\textbf{Partie II.}
\begin{enumerate}[resume]
\item Il revient à un sur dix (10.2~\% dans la simulation).
\item La latence passe de \qty{25.4}{\micro\second} à \qty{35.8}{\micro\second}.
\item Les copies ont chargé les passerelles aux moments de plus forte activité et rendu la latence plus élevée et moins prévisible pour chaque participant.
\item Un réseau déterministe, une file FIFO dans les passerelles, un point d'entrée unique propre à chaque partition~; des connexions de même longueur~; et
  des prix qui croissent avec le nombre de sessions.
\end{enumerate}
\textbf{Partie III.}
\begin{enumerate}[resume]
\item 9.8~\% avec l'une ou l'autre.
\item L'ordre dans lequel les messages atteignent la carte de la passerelle, et donc la vitesse propre de la firme vers la bourse, à ±2.5~ns de câble près.
\item La réaction propre de la firme~: traitement des flux de données de marché, décision et chemin jusqu'à la carte.
\item Supprimer l'avantage de découper un message en plusieurs paquets pour commencer à l'envoyer avant qu'il soit complet.
\end{enumerate}
\textbf{Partie IV.}
\begin{enumerate}[resume]
\item \emph{Résultat nommé}~: avec des passerelles parallèles et sujettes à la gigue, la chance qu'a une firme de gagner une course à dix firmes passe de 10.0~\% avec une
  seule session à 59.7~\% avec douze, alors que sous une passerelle de segment elle reste à 9.8~\%~; pour 20\,000 courses par mois et 0.50~euro la course, les
  onze sessions supplémentaires valent 4\,973~euros par mois, soit 452~euros par session en moyenne, mais chaque session au-delà de la sixième vaut moins
  que les 520~euros qu'elle coûte, et la douzième seulement 180.5~euros.
\item Que le tarif à partir de la septième est fixé là où la copie marginale vaut moins qu'elle ne coûte~: le barème impose ce que la
  conception n'a plus besoin.
\item Données publiques~: les conceptions des passerelles, les longueurs de chemin, les variances, l'égalité de longueur des câbles, les limites de session et les tarifs. Hypothèses~: les délais du modèle,
  les files d'attente, les temps de service, le nombre de firmes, les courses et leur valeur.
\item Parallèle~: des sessions jusqu'au point où leur valeur marginale rejoint le tarif, sachant que le gain disparaît si les autres copient. Segment~: des sessions
  pour le débit, les segments et la bascule (failover) seulement.
\item Mesurer~: envoyer des sondes horodatées à travers plusieurs sessions et passerelles, et comparer les horodatages de réception et d'appariement de la
  bourse~; si l'ordre au moteur diffère de l'ordre sur la \omterm{def:nw:networking-for-trading:frame}{carte réseau}, la conception n'est pas FIFO.
\item Jusqu'à $k$ exécutions par ordre à réconcilier par session, et des positions qui doivent supposer que toutes les copies s'exécutent.
\item Dans les lignes des contrats à terme~: segments, sessions et leurs tarifs, bascule, limites de compensation, drop copies, et les deux flux par des chemins distincts.
\item Que l'ordre d'arrivée à l'edge de la bourse, et non la chance d'un chemin à l'intérieur, décide qui arrive en premier.
\end{enumerate}
\end{solution}

\begin{solution}{iq:nw:futures-connectivity:1}
Une passerelle dédiée à un groupe de produits sur un même moteur, si bien que tous les ordres pour ceux-ci entrent par un seul chemin dans une seule séquence. Elle
supprime le routage à l'intérieur de la bourse et la variance des chemins parallèles, et donc la récompense liée à l'envoi de copies.
\iqlookfor{Segments~; une seule séquence~; équité~; routage supprimé.}
\end{solution}

\begin{solution}{iq:nw:futures-connectivity:2}
Des passerelles parallèles, avec des files indépendantes et de la gigue, font de l'arrivée de la firme le minimum des délais de ses copies. Elle cesse de fonctionner
lorsque la bourse ordonne les messages à un point d'entrée unique, lorsque tout le monde copie, ou lorsque les tarifs et les limitations de débit rendent les copies plus coûteuses que
ce qu'elles rapportent.
\iqlookfor{Statistiques d'ordre~; chemins indépendants~; équilibre~; réaction de la bourse.}
\end{solution}

\begin{solution}{iq:nw:futures-connectivity:3}
Envoyer des paires de messages à des instants connus par des sessions et des passerelles différentes, et comparer les horodatages de la bourse à son edge
réseau, avec l'ordre d'appariement~; mesurer la variance sur chaque chemin~; vérifier que des câbles égaux et des ports égaux donnent des temps égaux.
\iqlookfor{Sondes contrôlées~; horodatages de la bourse~; inversion d'ordre~; variance.}
\end{solution}

\begin{solution}{iq:nw:futures-connectivity:4}
Une cause commune~: le même commutateur, le même tampon, la même carte NIC ou le même hôte recevant les deux flux, ou une rafale supérieure à ce que cet élément partagé peut porter~; moins
probable, une perte à la bourse. Vérifier si les flux partagent un élément et les séparer.
\iqlookfor{Pertes corrélées~; éléments partagés~; reprise.}
\end{solution}

\begin{solution}{iq:nw:futures-connectivity:5}
Solder ou couvrir ce qui peut encore être négocié dans les actions autorisées, parler au desk de risque du membre compensateur de la limite et
du dépassement, et réconcilier les positions à partir du drop copy~; demain, dimensionner les stratégies sous la limite avec une marge, et demander des limites
par produit lorsqu'elles correspondent mieux au risque.
\iqlookfor{Risque immédiat~; communication~; réconciliation~; limites dans le plan.}
\end{solution}

\begin{solution}{iq:nw:futures-connectivity:6}
Une seule identité de session selon le modèle de la bourse, connectée à la passerelle de chaque segment, principale et de secours, avec des limitations de débit dimensionnées par
segment~; pas de copies avec une conception FIFO~; bascule testée pendant les fenêtres de la bourse~; limites de compensation connues et suivies~; drop copies
réconciliées par session~; les deux flux de données de marché sur des chemins distincts.
\iqlookfor{Segments~; débit~; bascule~; risque~; réconciliation~; flux.}
\end{solution}
